%-------------------------------------------------------------------------------
%	WORK PLAN & REVIEW PAGE
%-------------------------------------------------------------------------------

%---------------------ТӨЛӨВЛӨГӨӨНИЙ ХУУДСЫН ЭХЛЭЛ--------------------------------
\begin{titlepage}

\vspace*{0.5cm}
Батлав. \deptname йн эрхлэгч: 
\begin{flushright}
\makebox[4cm]{\dotfill} /\chairname/ 
\end{flushright}

Удирдагч: 
\begin{flushright}
\makebox[4cm]{\dotfill} /\supname/
\end{flushright}

\begin{center}

\vspace*{2cm}
\textbf{{\large ТӨГСӨЛТИЙН АЖЛЫН \\ ТӨЛӨВЛӨГӨӨ}}\\[0.5cm]

\textsc{\large Сэдэв: ''\ttitle''}\\[0.5cm]

\begin{tabular}{|c|p{7cm}|c|c|}
	%\rowcolor{gray!40}
	\hline
	№ & \makebox[7cm][c]{Судалгааны ажлын бүлэг, хэсгүүдийн нэр} & Эзлэх хувь & Гүйцэтгэх хугацаа \\ \hline
	1 & {1. Онолын судалгаа} & & \\
	& 1.1 Веб сайт гэж юу вэ? & & \\
	& 1.2 HTML хэл, бүтэц & & \\ 
	& 1.3 CSS гэж юу вэ? &  & \\ 
	& 1.4 JavaScript гэж юу вэ? & 30\% &  2023.03.05\\
	& 1.4 React програмчлалын хэл & & \\
    & 1.5 Өгөгдлийн сан гэж юу вэ? & & \\
    & 1.6 MongoDB гэж юу вэ? & & \\ \hline
    
	2 & {2. Системийн шаардлага тодорхойлох, шинжилгээ хийх} & & \\
	& 2.1 Системийг ашиглах хэрэглэгчид & & \\
    & 2.2 Функционал шаардлага & 30\% & 2023.03.20\\
    & 2.3 Функционал бус шаардлага& & \\
    & 2.4 Юзкейс диаграм& & \\
    & 2.5 Юзкейс диаграмын тодорхойлолт& & \\ \hline

	3 & {3. Системийн зохиомж}     &   &  \\
    & 3.1 Дэлгэцийн зохиомж & 30\%&2023.04.15 \\
    & 3.2 Үйл ажиллагааны диаграммууд & &\\ \hline
	4 & {Ерөнхий дүгнэлт} & 10\% & 2023.05.24 \\ \hline
\end{tabular}

\vspace{2cm}
Төлөвлөгөөг боловсруулсан оюутан: \makebox[3cm]{\dotfill} /\shortname/\\

\vspace{1cm}
Төлөвлөгөөг хянасан: \makebox[3cm]{\dotfill} /\chairname/

\end{center}
\end{titlepage}
%---------------------ТӨЛӨВЛӨГӨӨНИЙ ХУУДСЫН ТӨГСГӨЛ ---------------------------------------------------


%---------------------ГҮЙЦЭТГЭЛИЙН ХУУДСЫН ЭХЛЭЛ ---------------------------------------------------
\newpage

\begin{titlepage}

\begin{center}

\vspace*{2cm}
\textbf{{\large ТӨГСӨЛТИЙН АЖЛЫН ЯВЦ}}\\[0.5cm]

\begin{tabular}{|c|p{7cm}|c|c|}
	%\rowcolor{gray!40}
	\hline
	№ & \makebox[7cm][c]{Төслийн бүлэг хэсгүүдийн нэр} & Биелсэн     & Удирдагчийн \\
	  &                    & хугацаа    & гарын үсэг \\ \hline
	1 & {1. Онолын судалгаа} & & \\
	& 1.1 Веб сайт гэж юу вэ? & & \\
	& 1.2 HTML хэл, бүтэц & & \\ 
	& 1.3 CSS гэж юу вэ? &  & \\ 
	& 1.4 JavaScript гэж юу вэ? & III.05 & \\
	& 1.4 React програмчлалын хэл & & \\
    & 1.5 Өгөгдлийн сан гэж юу вэ? & & \\
    & 1.6 MongoDB гэж юу вэ? & & \\ \hline
    
	2 & {2. Системийн шаардлага тодорхойлох, шинжилгээ хийх} & & \\
	& 2.1 Системийг ашиглах хэрэглэгчид & & \\
    & 2.2 Функционал шаардлага & III.20 & \\
    & 2.3 Функционал бус шаардлага& & \\
    & 2.4 Юзкейс диаграм& & \\
    & 2.5 Юзкейс диаграмын тодорхойлолт& &  \\ \hline
    
	3 & {3. Системийн зохиомж}     &  & \\
    & 3.1 Дэлгэцийн зохиомж &IV.15  & \\
    & 3.2 Үйл ажиллагааны диаграммууд & &\\ \hline

    4 & {Ерөнхий дүгнэлт} & V.24 & \\ \hline
\end{tabular}


\end{center}

\begin{center}
Удирдагч: \makebox[3cm]{\dotfill} /\supname/ \\

\vspace{5cm}
ЗӨВШӨӨРӨЛ \\[0.5cm]
Оюутан \shortname --н бичсэн төгсөлтийн ажлыг УШК-д хамгаалуулахаар тодорхойлов.\\[0.5cm]
Салбарын эрхлэгч: \makebox[3cm]{\dotfill} /\chairname/
\end{center}

\end{titlepage}
%---------------------ГҮЙЦЭТГЭЛИЙН ХУУДСЫН ТӨГСГӨЛ ---------------------------------------------------